\chapterbackmatter{Weinberg Exercises}

\begin{enumerate}
\WeinbergExercise{1} Calculate the effective ghost Lagrangian in a `generalized unitarity
gauge,' with $B[f]$ given as usual by
Eq.~\hyperref[eq:15.5.22]{(15.5.22)}, but now where
$f_\alpha(x)=i\phi_n(x)(t_\alpha)_{nm}
\langle\phi_m(0)\rangle_{\rm VAC}$, with $\phi_n$ real scalar fields, and
$t_\alpha$ imaginary antisymmetric matrices representing the Lie algebra of
the gauge group. What is the ghost propagator? Is this part of the
Lagrangian renormalizable?

\WeinbergExercise{2} What would be the effect in the $SU(2)\cdot U(1)$ electroweak theory if
the gauge symmetry were broken by the vacuum expectation value of a field
$\phi_3$ belonging to a real triplet
$\vec\phi=(\phi^+,\phi^0,\phi^-)$ instead of the usual complex doublet
$(\phi^0,\phi^-)$?

\WeinbergExercise{3} Consider the standard electroweak theory, with a single scalar doublet.
In one-loop order, calculate the effect of $Z^0$ and neutral scalar exchange
on the anomalous magnetic moment of the muon.

\WeinbergExercise{4} What is the lowest-order magnetic moment of the $W^+$ and $Z^0$
particles in the standard electroweak theory?

\WeinbergExercise{5} What would be the effect of the discovery of a fourth generation of
quarks and leptons on the predictions of the unified theories of strong and
electroweak interactions discussed in
\hyperref[sec:21.5]{Section 21.5}?

\WeinbergExercise{6} Suppose that several fields with incommensurate values of the electric
charge had nonvanishing vacuum expectation values in a superconductor. What
effect would this have on the properties of superconductors discussed in
\hyperref[sec:21.6]{Section 21.6}?
\end{enumerate}
